% ============================================
% Macroeconomic Research Associate Cover Letter
% Dodge & Cox
% ============================================

\documentclass[11pt,letterpaper]{letter}

\usepackage{geometry}
\usepackage{microtype}
\usepackage[hidelinks]{hyperref}

\geometry{left=25mm,right=25mm,top=25mm,bottom=25mm}

\signature{Decory Edwards}
\address{
Department of Economics \\
Johns Hopkins University \\
Baltimore, MD 21218 \\
\texttt{dedwar65@jh.edu}
}

\begin{document}

\begin{letter}{
Hiring Committee \\
Dodge \& Cox \\
San Francisco, CA
}

\opening{Dear Members of the Hiring Committee,}

I am writing to apply for the 2026 Macroeconomic Research Associate position at Dodge \& Cox. I am a Ph.D.\ candidate in Economics at Johns Hopkins University, advised by Professors Christopher D.~Carroll, Nicholas W.~Papageorge, and Stelios Fourakis, and I expect to complete my degree in May 2026. My training is in macroeconomics, quantitative analysis, and applied data work, with a particular focus on long-run economic dynamics, heterogeneity, and policy-relevant macroeconomic risk.

My research agenda focuses on understanding how heterogeneity in economic environments shapes aggregate outcomes over long horizons. In my job market paper, I estimate persistent differences in individual returns to wealth using micro data from the Survey of Consumer Finances and embed these estimates in a structural macroeconomic model. I show that allowing for heterogeneity in returns materially improves the model’s ability to match observed wealth concentration and has important implications for aggregate saving behavior, capital accumulation, and the transmission of macroeconomic shocks. While developed in an academic setting, this framework reflects the same core challenge faced by global investors, namely assessing how persistent structural differences translate into long-run macroeconomic risk and opportunity.

In a related project using the Health and Retirement Study, I examine how beliefs and interpersonal trust affect portfolio performance and realized returns. I document a hump-shaped relationship between trust and returns to wealth, highlighting a behavioral channel that links expectations, risk-taking, and realized outcomes. This work reinforces my broader interest in understanding how beliefs, institutions, and economic structure interact to shape asset returns over time.

I am particularly drawn to Dodge \& Cox because of its long-term, fundamentals-driven investment philosophy and its emphasis on rigorous macroeconomic analysis across countries, asset classes, and currencies. The opportunity to conduct in-depth country research, monitor global economic and financial data, and contribute to interest rate and currency analysis aligns closely with my training. I have extensive experience working with large macroeconomic datasets, building and maintaining analytical pipelines in Python, and translating complex quantitative results into clear economic narratives. I am especially excited by the chance to support sovereign and global fixed income investment decisions through careful assessment of structural trends, vulnerabilities, and policy regimes.

The collaborative nature of the Macroeconomic Research Associate role is also very appealing to me. I value environments where research directly informs decision making, and where quantitative analysis is combined with judgment and institutional knowledge. I am comfortable working across teams, responding to internal research requests, and contributing to investment and client-facing materials. I am also enthusiastic about Dodge \& Cox’s hybrid work structure and the opportunity to be fully engaged with colleagues in the San Francisco office.

I would welcome the opportunity to contribute my macroeconomic training, quantitative skills, and research experience to Dodge \& Cox’s investment process. Thank you very much for your time and consideration.

\closing{Sincerely,}

\end{letter}
\end{document}
